% Fragmento público generado desde propietarios íntegros.
% Residencia: 52.2.2.
% Fuente: ../incoming/radion_monografia_20260827/manuscrito/sections/03_cierre_exacto.tex:51-103

\paragraph{Densité orientée \(20/81\) du secteur actif \(90/120\)}

Le sélecteur tritique de phase sépare
\[
H_+=\{1,4,7\},\qquad H_-=\{2,5,8\},\qquad H_0=\{3,6,9\}.
\]
Il y a six phases actives et quinze paires non ordonnées :
\[
H_{\mathrm{act}}=H_+\cup H_-,\qquad
\left|\binom{H_{\mathrm{act}}}2\right|=\binom62=15.
\]
Le mode actif et son extension à huit positions sont
\begin{equation}
N_6=6\binom62=90,\qquad
N_8=8\binom62=120.
\label{eq:90-120-incidencia}
\end{equation}
Deux feuillets conjugués du secteur actif, normalisés par la fibre ambiante,
donnent
\begin{equation}
\rho_{\mathrm{tw}}
=\frac{2N_6}{|\F_3^6|}
=\frac{2\cdot90}{729}
=\frac{180}{729}
=\boxed{\frac{20}{81}}.
\label{eq:rho-twoway}
\end{equation}

\begin{proposition}[Unicité relative de la densité]
Le secteur actif \(N_6=90\), les deux feuillets conjugués et la
normalisation par \(|\F_3^6|=729\) étant fixés, la densité orientée est nécessairement
\(20/81\).
\end{proposition}
\begin{proof}
La cardinalité transportée est \(2N_6=180\). La normalisation par l'espace ambiant
donne \(180/729\) ; diviser le numérateur et le dénominateur par \(9\) produit \(20/81\).
Aucune comparaison avec \(180/\pi\) ni avec \(\epsR\) n'intervient.
\end{proof}

La lecture historique comme projection sur un réseau \(1/81\) est un contrôle de
rigidité utile, mais n'est plus le fondement : le coefficient s'obtient avant
la clôture. L'ablation d'un feuillet produit \(10/81\), et la sortie change à
l'ordre \(10^{-5}\) ; le facteur deux n'est donc pas décoratif.

La section torsionnelle est
\begin{equation}
\Phi_T=\rho_{\mathrm{tw}}\Delta_4,
\qquad
S_T=1+\Phi_T.
\label{eq:seccion-torsional}
\end{equation}

